\begin{remarks}[7.5.4]
\EGAReaderTarget{ega4.c04.item.7-5-4}{IV.2.s7.5.4-fr}
\begin{enumerate}
\item[(i)] Sería interesante saber si el teorema (7.5.1) sigue
siendo válido sin la hipótesis de finitud (i) sobre el cuerpo residual
de $B$. La respuesta es afirmativa cuando el cuerpo residual $k$ de
$A$ tiene característica 0, como se ve utilizando los resultados de
Hironaka sobre la resolución de singularidades (cf. (7.9.8)).
Señalemos el siguiente caso particular de la cuestión planteada aquí:
sean $A$ y $B$ dos anillos noetherianos completos, $k$ el cuerpo
residual de $A$ y $\varphi:A\to B$ un homomorfismo local que hace de
$B$ una $A$-álgebra \emph{formalmente suave} (\textbf{0}, 19.3.1) (lo que equivale a
decir que $B$ es un $A$-módulo plano y que $B\otimes_A k$ es
geométricamente regular sobre $k$ (\textbf{0}, 19.7.1)). Entonces, ¿son
geométricamente regulares las fibras del morfismo
${}^a\varphi:\operatorname{Spec}(B)\to\operatorname{Spec}(A)$?
Así ocurre cuando el cuerpo residual $k'$ de $B$ es una extensión
\emph{finita} de $k$, según (7.5.1); se puede demostrar que la respuesta
sigue siendo afirmativa cuando $k'$ es una extensión \emph{de tipo finito}
de $k$. Ignoramos la respuesta cuando $B\otimes_A k$ es un cuerpo
igual a la \emph{clausura separable} de $k$.

\item[(ii)] Podría enunciarse un resultado análogo a (7.5.1)
relativo a un $A$-módulo $M$ y un $B$-módulo $N$, ambos de tipo finito,
sobre las propiedades de
$(M\otimes_A N)_y\otimes_{\mathscr O_x}\boldsymbol{k}(y)$,
donde $y$ recorre $\operatorname{Spec}(B)$ y $x=f(y)$, que se deducen
de propiedades del mismo tipo de $M_x$ y $N_y$.

\item[(iii)] Si la propiedad $\boldsymbol R$ satisface las
condiciones $(\mathrm R'_{\mathrm I})$, $(\mathrm R_{\mathrm{II}})$,
$(\mathrm R_{\mathrm{III}})$ y $(\mathrm R_{\mathrm{IV}})$, y se
cumplen las hipótesis (i) y (ii) de (7.5.2), entonces, de las
propiedades $\boldsymbol R(A)$ y
$\boldsymbol P(\operatorname{Spec}(\widehat B
\otimes_{\widehat A}k),k)$
se deduce $\boldsymbol R(B)$. Como se ha visto en (7.5.3), así ocurre para
las propiedades dadas como ejemplos en (7.5.3), salvo para (i) y (vii).
En cuanto a (vii), parece, sin embargo, plausible que la respuesta a la
siguiente pregunta sea afirmativa: sea $\varphi:A\to B$ un
homomorfismo local de anillos locales noetherianos que hace de $B$ un
$A$-módulo \emph{plano}. Supongamos que $A$ es \emph{completo}, \emph{íntegro} y
\emph{geométricamente unirrama}, y que la fibra
$\operatorname{Spec}(B\otimes_A k)$ (donde $k$ es el cuerpo residual
de $A$) es \emph{geométricamente localmente íntegra}; ¿es entonces cierto
que $B$ es \emph{íntegro}? Esto puede demostrarse cuando $k$ tiene
característica 0, utilizando la resolución de singularidades de
Hironaka (7.9.8). También puede mostrarse que la respuesta es
afirmativa cuando $B$ es una $A$-álgebra esencialmente de tipo finito
(1.3.8) (cf. (11.3.10) y (11.3.11)) o cuando
$\operatorname{Spec}(B\otimes_A k)$ es \emph{geométricamente normal}
(pues de (7.5.3) resulta entonces que el morfismo
$\operatorname{Spec}(B)\to\operatorname{Spec}(A)$ es normal, y se
concluye por (6.15.10)). Pero, aun suponiendo que el cuerpo residual de
$B$ es igual al de $A$ y que $B$ también es completo, ignoramos si la
respuesta es afirmativa en el caso general.

\item[(iv)] Consideremos la propiedad $\boldsymbol R(C)$:
«$C$ es reducido, equidimensional y satisface $(\mathrm R_n)$».
Satisface $(\mathrm R_{\mathrm{IV}})$, en virtud de (5.12.5), pero
se ignora, en cambio, si satisface $(\mathrm R'_{\mathrm I})$,
$(\mathrm R_{\mathrm{II}})$ o $(\mathrm R_{\mathrm{III}})$; las
dificultades proceden de la verificación de la equidimensionalidad.
\end{enumerate}
\end{remarks}
